






\documentclass{tietreport}

% Prevent fancyhdr warning about insufficient header height
\setlength{\headheight}{14.5pt}

% Import images
\usepackage{graphicx}
\graphicspath{{figures/}{./}}

% Tables and maths
\usepackage{amsmath}
\usepackage{array}
\usepackage{tabularx}
% Software-design diagrams

% Bibliography configuration
\usepackage[
  date=year,
  style=alphabetic,
  backend=biber,
  sorting=nty,
  maxnames=5,
  minnames=3,
  maxalphanames=4,
  minalphanames=3,
  backref=true,
  doi=false,
  isbn=false,
  url=true,
  eprint=false
]{biblatex}
\addbibresource{references.bib}

%% ==================================================
%% PROJECT DETAILS
%% ==================================================

\TitlePageHeader{%
  UCS503P%
}

\ProjectTitle{%
  Personal Finance Manager%
}

\ProjectSubTitle{%
  A Desktop-Based Financial Tracking and Budgeting System%
}

\author{%
  \texttt{1024030022} Arjun Puri \\
  \texttt{1024030105} Anshdeep Singh%
}

\TitlePageSubText{%
  Group: \texttt{3C11} BE Third Year, CSE%
}

\AdvisorName{%
  Ms. Parmveer%
}

\TitlePageFooterText{{%
    \large Final Project Report \\[0.2cm]
    \bfseries Computer Science and Engineering
    Department%
  } \\[0.15cm]
  Thapar Institute of Engineering and Technology,
  Patiala%
}

%% ==================================================

\begin{document}
\raggedbottom

% Removes the blank pages caused by chapters opening on odd pages.
% Delete this line if your department requires chapters on right-hand pages.
\let\cleardoublepage\clearpage

\maketitle

\tableofcontents
\listoffigures
\listoftables

% ===== CHAPTER 1: INTRODUCTION =====

\chapter{Introduction}

\section{Background and Context}

Students and young professionals usually earn from several
sources, such as an allowance, a stipend or part-time work, and
spend across many categories, including food, travel,
subscriptions, shopping and bills. Without a structured system it
is hard to see where money is going or whether monthly spending is
under control. The \emph{Personal Finance Manager} (PFM) is a
desktop application that brings income, expenses, budgets and
savings goals into one place and turns them into summaries that
the user can read at a glance.

The project was proposed in~\cite{pfmproposal} with a
\emph{time-to-value} approach: build the login and
transaction-and-database core first, validate it, and then layer
budgets, goals and reports on top. This report documents the
implemented desktop application, its design, and the results of
the first phase.

\subsection{Problem Statement}

Tracking personal finances by hand or across scattered sources
leads to four recurring problems:

\begin{itemize}
  \item \textbf{Fragmentation:} transactions are spread across
    bank alerts, UPI apps and cash, with no single source of truth.
  \item \textbf{Low visibility:} there is no immediate view of how
    actual spending compares with a monthly budget.
  \item \textbf{No goal tracking:} savings targets, such as a
    laptop or a phone, are rarely tracked with a visible progress
    measure.
  \item \textbf{Manual, repetitive entry:} recurring costs such as
    rent are re-entered every month.
\end{itemize}

In its current version PFM provides a single MySQL-backed store for
records the user enters, which partly addresses fragmentation; it
does not import data from banks or UPI apps. It addresses low
visibility and goal tracking directly. Manual entry is not
reduced: recurring-expense handling is planned for a later
iteration (Section~\ref{sec:limitations}).

\section{Project Objectives}

\begin{enumerate}
  \item Provide a login and a single MySQL-backed store for a
    user's income and expense records, with validation of amount,
    date and required fields.
  \item Show the user's financial position (income, expenses, net
    savings, category breakdown and recent activity) on a
    dashboard that updates from the stored data.
  \item Let the user set monthly spending limits per category and
    see usage against each limit.
  \item Let the user create savings goals with a target, amount
    saved and optional deadline, and view progress towards each
    goal.
  \item Provide monthly reports with a category breakdown and a
    CSV export.
  \item Keep the presentation, logic and data layers separate so
    the system can later move to a web back end without
    redesigning the schema.
\end{enumerate}

Objectives~1--5 are implemented in the current build, although the
login does not yet hash passwords (Section~\ref{sec:limitations}).
Objective~6 is partly met and is discussed under future work.

% ===== CHAPTER 2: METHODOLOGY & DESIGN =====

\chapter{Methodology and Design}

\section{System Architecture}

\subsection{High-Level Design}

PFM is designed around the three-tier structure set out in the
proposal. The tiers are only partly separated in the current
build (Section~\ref{sec:limitations}).

\begin{itemize}
  \item \textbf{Presentation tier:} a Tkinter GUI with a login
    screen and eight pages reachable from a persistent sidebar
    (Dashboard, Transactions, Income, Expenses, Budget, Savings
    Goals, Reports and Settings).
  \item \textbf{Business-logic tier:} Python code that validates
    input and computes the totals, percentages and budget figures
    shown on screen. Some of this code still lives inside the GUI
    classes.
  \item \textbf{Data tier:} a MySQL database accessed through
    PyMySQL, with queries parameterised using placeholders.
\end{itemize}

The core workflow is: the user logs in; the dashboard loads;
the user adds an income or expense record; the input is
validated; the record is written to MySQL; and the totals,
budget usage and charts are recalculated when the views refresh.

\subsection{Navigation Architecture}

The application uses a single controller window (the \texttt{App}
class, derived from \texttt{tk.Tk}) that owns all window-level
settings. Each page is a \texttt{tk.Frame} built once and stacked
in a shared container with \texttt{grid()}. Switching pages calls
\texttt{show\_frame()}, which raises the chosen frame with
\texttt{tkraise()}, and each page exposes an \texttt{on\_show()}
hook that reloads its data when the page becomes visible. This
replaced an earlier design that destroyed and recreated widgets on
every navigation, which was slower and harder to maintain.

\subsection{Technical Stack}

\begin{itemize}
  \item \textbf{Language:} Python 3, using the standard library
    wherever possible.
  \item \textbf{GUI:} Tkinter~\cite{pythontkinter}, with the
    \texttt{tkcalendar} \texttt{DateEntry} widget for date
    selection~\cite{tkcalendardocs}.
  \item \textbf{Database:} MySQL~\cite{mysqlmanual} accessed
    through PyMySQL~\cite{pymysqldocs}.
  \item \textbf{Charts:} drawn directly on \texttt{tk.Canvas}
    (bar, line, donut and progress-bar views), with no external
    plotting library.
  \item \textbf{Export:} CSV export of the monthly report.
\end{itemize}

\subsection{Financial Calculations}

All displayed figures are deterministic formulas, so every number
can be audited by hand:

\begin{align*}
  \text{Net Savings} &= \text{Total Income} - \text{Total Expenses}\\
  \text{Budget Usage (\%)} &= \frac{\text{Amount Spent}}{\text{Budget Limit}} \times 100\\
  \text{Goal Progress (\%)} &= \frac{\text{Amount Saved}}{\text{Target Amount}} \times 100\\
  \text{Category Share (\%)} &= \frac{\text{Category Expense}}{\text{Total Expenses}} \times 100
\end{align*}

``Amount Spent'' is not defined identically everywhere. On the
Budget page it is the expense total over categories that have a
budget; on the sidebar spending-limit card it is the total of all
expenses. Percentages are shown as whole numbers, and the rounding
is not yet identical on every page. Both points are discussed in
Section~\ref{sec:analysis}.

\section{Software Design Models}

The following diagrams describe the implemented user-facing workflows and the
logical design of PFM. The database diagram is a conceptual model derived from
the features documented in this report. Attribute names and relationships
should be checked against the final MySQL \texttt{SHOW CREATE TABLE} output
before treating it as the physical schema.

\subsection{Use-Case Diagram}

Figure~\ref{fig:usecase} presents the main interactions available to a
logged-in user. Authentication is shown separately because the application
requires a successful login before the finance pages can be accessed.

\begin{figure}[htbp]
\centering
\begin{center}
\fbox{\begin{minipage}{0.92\textwidth}
\centering
\textbf{PERSONAL FINANCE MANAGER -- USER INTERACTIONS}

\vspace{0.25cm}
\begin{tabular}{c|c}
\textbf{External actor} & \textbf{System use cases}\\
\hline
\textbf{User} $\longrightarrow$ & Log in\\
 & View dashboard\\
 & View transactions\\
 & Add income\\
 & Add / delete expenses\\
 & Manage budgets\\
 & Manage savings goals\\
 & Generate monthly report / export CSV\\
 & View account / change password
\end{tabular}
\end{minipage}}
\end{center}
\caption{Use-case diagram of the Personal Finance Manager}
\label{fig:usecase}
\end{figure}

\subsection{System Architecture Diagram}

The architecture follows a presentation, business-logic and data-access
separation. In the current implementation, some calculations and database
queries remain inside GUI page classes; therefore, the diagram represents the
logical responsibilities rather than a claim that all layers are already
independent.

\begin{figure}[htbp]
\centering
\begin{center}
\fbox{\begin{minipage}{0.90\textwidth}
\centering
\textbf{LOGICAL THREE-TIER SYSTEM ARCHITECTURE}

\vspace{0.2cm}
\fbox{\parbox{0.78\textwidth}{\centering \textbf{User}}}
\\[0.12cm] $\Downarrow$
\\[0.12cm]
\fbox{\parbox{0.78\textwidth}{\centering \textbf{Presentation tier}\\Tkinter application and App controller}}
\\[0.12cm] $\Downarrow$
\\[0.12cm]
\begin{tabular}{|c|c|}
\hline
Dashboard, Transactions & Budget, Savings Goals\\
Income, Expenses & Reports, Settings\\
\hline
\end{tabular}
\\[0.12cm] $\Downarrow$
\\[0.12cm]
\fbox{\parbox{0.78\textwidth}{\centering \textbf{Business logic}\\Input validation, totals, percentages, budget and goal calculations}}
\\[0.12cm] $\Downarrow$
\\[0.12cm]
\fbox{\parbox{0.78\textwidth}{\centering \textbf{Data access}\\PyMySQL parameterised queries}}
\\[0.12cm] $\Downarrow$
\\[0.12cm]
\fbox{\parbox{0.78\textwidth}{\centering \textbf{MySQL database}\\Users, income, expenses, budgets, savings goals and related records}}
\\[0.12cm]
\textit{Query results return from the database to the application for display.}
\end{minipage}}
\end{center}
\caption{Logical three-tier architecture of PFM}
\label{fig:architecture}
\end{figure}

\subsection{Activity Diagram: Recording an Expense}

The activity diagram in Figure~\ref{fig:activity} describes the expense-entry
workflow, including validation and persistence. The same general validation
pattern is used for income entry.

\begin{figure}[htbp]
\centering
\begin{center}
\fbox{\begin{minipage}{0.90\textwidth}
\centering
\textbf{EXPENSE-ENTRY ACTIVITY FLOW}

\vspace{0.2cm}
\fbox{\parbox{0.68\textwidth}{\centering Start}}
\\[0.10cm] $\Downarrow$
\\[0.10cm]
\fbox{\parbox{0.68\textwidth}{\centering Open Expenses page}}
\\[0.10cm] $\Downarrow$
\\[0.10cm]
\fbox{\parbox{0.78\textwidth}{\centering Enter expense name, amount, category and date}}
\\[0.10cm] $\Downarrow$
\\[0.10cm]
\fbox{\parbox{0.78\textwidth}{\centering Validate required fields, amount and date}}
\\[0.10cm] $\Downarrow$
\\[0.10cm]
\fbox{\parbox{0.55\textwidth}{\centering \textbf{Is the input valid?}}}
\\[0.15cm]
\begin{tabular}{|c|c|}
\hline
\textbf{No} & \textbf{Yes}\\
\hline
Show warning and return & Save expense in MySQL\\
to correct the input & $\Downarrow$\\
 & Refresh expense list\\
 & and financial summaries\\
\hline
\end{tabular}
\\[0.12cm] $\Downarrow$
\\[0.08cm] \fbox{\parbox{0.68\textwidth}{\centering End}}
\end{minipage}}
\end{center}
\caption{Activity diagram for adding an expense}
\label{fig:activity}
\end{figure}

\subsection{Sequence Diagram: Saving a Financial Record}

Figure~\ref{fig:sequence} shows the typical interaction between the user,
Tkinter page, validation and database during a record-creation operation.
The database response is followed by a refresh so that the visible totals and
transaction lists reflect persisted data.

\begin{figure}[htbp]
\centering
\begin{center}
\fbox{\begin{minipage}{0.96\textwidth}
\centering
\textbf{SEQUENCE: SAVING A FINANCIAL RECORD}

\vspace{0.2cm}
\begin{tabular}{|p{2.4cm}|p{2.4cm}|p{2.6cm}|p{2.0cm}|}
\hline
\textbf{User} & \textbf{Tkinter page} & \textbf{Validation / logic} & \textbf{MySQL}\\
\hline
Enter record and submit & $\leftarrow$ Receive form data & &\\
\hline
 & Validate fields & $\rightarrow$ Check values &\\
\hline
 & $\leftarrow$ Validation result & Return valid / invalid &\\
\hline
 & Parameterised INSERT & & $\rightarrow$ Store record\\
\hline
 & $\leftarrow$ Show success / error & & Return query result\\
\hline
 & Reload records and totals & & $\rightarrow$ Read updated data\\
\hline
View refreshed page & $\leftarrow$ Display updated data & & Return records\\
\hline
\end{tabular}
\end{minipage}}
\end{center}
\caption{Sequence diagram for saving and displaying a financial record}
\label{fig:sequence}
\end{figure}

\subsection{Entity--Relationship Diagram}

The ER model in Figure~\ref{fig:er} is derived from the actual
\texttt{CREATE TABLE} statements in the supplied MySQL dump for
\texttt{personal\_finance\_manager}. It distinguishes enforced foreign-key
relationships (solid arrows) from fields that appear to associate records
at application level but are not protected by a declared foreign key
(dashed arrows). This distinction is important because a matching column
name does not by itself establish referential integrity.

\begin{figure}[htbp]
\centering
\begin{center}
\small
\setlength{\tabcolsep}{4pt}
\renewcommand{\arraystretch}{1.25}
\begin{tabularx}{\textwidth}{|>{\raggedright\arraybackslash}p{0.19\textwidth}|X|>{\raggedright\arraybackslash}p{0.27\textwidth}|}
\hline
\textbf{Entity} & \textbf{Primary key and important attributes} & \textbf{Declared relationships}\\
\hline
\texttt{users} & PK: \texttt{user\_id}; username, email, password & Parent of categories, transactions, savings goals and recurring expenses\\
\hline
\texttt{categories} & PK: \texttt{category\_id}; FK: \texttt{user\_id}; name, type & N:1 to users; parent of recurring expenses\\
\hline
\texttt{transactions} & PK: \texttt{transaction\_id}; FK: \texttt{user\_id}; category\_id, amount, type, txn\_date & N:1 to users; category\_id has no declared FK\\
\hline
\texttt{savings\_goals} & PK: \texttt{goal\_id}; FK: \texttt{user\_id}; name, target/current amount & N:1 to users; parent of goal contributions\\
\hline
\texttt{goal\_contributions} & PK: \texttt{contribution\_id}; FK: \texttt{goal\_id}; amount, source, created\_at & N:1 to savings goals\\
\hline
\texttt{recurring\_expenses} & PK: \texttt{recurring\_id}; FKs: \texttt{user\_id}, \texttt{category\_id}; amount, frequency, next due date & N:1 to users and categories\\
\hline
\texttt{income} & PK: \texttt{id}; username, amount, category, income date, source & Username/category associations have no declared FK\\
\hline
\texttt{expenses} & PK: \texttt{id}; username, name, amount, category, expense date & Username/category associations have no declared FK\\
\hline
\texttt{budgets} & PK: \texttt{id}; username, category, monthly limit & Username/category associations have no declared FK\\
\hline
\texttt{reports} & PK: \texttt{report\_id}; user\_id, report name, month/year and totals & \texttt{user\_id} has no declared FK\\
\hline
\end{tabularx}
\end{center}
\caption{Entity and relationship summary based on the supplied MySQL schema.
Declared foreign keys are distinguished from application-level associations.}
\label{fig:er}
\end{figure}

\subsection{Relational Schema Summary}

Table~\ref{tab:dbschema} summarises the ten tables and their key columns.
The complete SQL dump defines additional data types, defaults, indexes,
unique constraints and delete rules; the table below is a compact
documentation view of the physical schema.

\begin{table}[htbp]
\centering
\scriptsize
\begin{tabularx}{\textwidth}{|p{2.5cm}|p{2.5cm}|X|p{3.0cm}|}
\hline
\textbf{Table} & \textbf{Primary key} & \textbf{Important columns} & \textbf{Declared foreign key(s)}\\
\hline
\texttt{users} & \texttt{user\_id} & \texttt{username}, \texttt{email}, \texttt{password} & None\\
\hline
\texttt{categories} & \texttt{category\_id} & \texttt{name}, \texttt{type}; unique per user/name/type & \texttt{user\_id} $\rightarrow$ \texttt{users.user\_id}\\
\hline
\texttt{income} & \texttt{id} & \texttt{username}, \texttt{amount}, \texttt{category}, \texttt{income\_date}, \texttt{source} & None\\
\hline
\texttt{expenses} & \texttt{id} & \texttt{username}, \texttt{name}, \texttt{amount}, \texttt{category}, \texttt{expense\_date} & None\\
\hline
\texttt{transactions} & \texttt{transaction\_id} & \texttt{user\_id}, \texttt{category\_id}, \texttt{amount}, \texttt{type}, \texttt{txn\_date} & \texttt{user\_id} $\rightarrow$ \texttt{users.user\_id}\\
\hline
\texttt{budgets} & \texttt{id} & \texttt{username}, \texttt{category}, \texttt{monthly\_limit} & None\\
\hline
\texttt{savings\_goals} & \texttt{goal\_id} & \texttt{user\_id}, \texttt{name}, \texttt{target\_amount}, \texttt{current\_amount}, \texttt{auto\_percent} & \texttt{user\_id} $\rightarrow$ \texttt{users.user\_id}\\
\hline
\texttt{goal\_contributions} & \texttt{contribution\_id} & \texttt{goal\_id}, \texttt{amount}, \texttt{note}, \texttt{source} & \texttt{goal\_id} $\rightarrow$ \texttt{savings\_goals.goal\_id}\\
\hline
\texttt{recurring\_expenses} & \texttt{recurring\_id} & \texttt{user\_id}, \texttt{category\_id}, \texttt{amount}, \texttt{frequency}, \texttt{next\_due\_date}, \texttt{is\_active} & \texttt{user\_id} $\rightarrow$ \texttt{users.user\_id}; \texttt{category\_id} $\rightarrow$ \texttt{categories.category\_id}\\
\hline
\texttt{reports} & \texttt{report\_id} & \texttt{user\_id}, \texttt{report\_name}, \texttt{report\_month}, \texttt{report\_year}, totals & None\\
\hline
\end{tabularx}
\caption{Summary of the physical MySQL tables and key relationships}
\label{tab:dbschema}
\end{table}

\subsection{Referential Integrity Observations}

The supplied SQL schema declares foreign keys for categories, transactions,
savings goals, goal contributions and recurring expenses. In particular,
\texttt{goal\_contributions.goal\_id} references
\texttt{savings\_goals.goal\_id} with cascading deletion, and the user-owned
category, transaction, savings-goal and recurring-expense records reference
\texttt{users.user\_id} as defined in the dump.

Some other columns resemble user or category references but do not have
foreign-key constraints in the supplied schema. These include
\texttt{income.username}, \texttt{expenses.username},
\texttt{budgets.username}, \texttt{reports.user\_id} and
\texttt{transactions.category\_id}. The latter has an index but no declared
foreign key to \texttt{categories}. Consequently, the database itself does
not enforce those associations. They may be maintained by application
queries, but this should not be described as database-enforced referential
integrity.

The \texttt{reports} table stores monthly report metadata and totals, while
the report screen described earlier can also calculate a report from
transaction records. The SQL definition alone does not establish whether
every generated report is persisted in this table.
\subsection{Data-Flow Diagram}

At context level, the user supplies authentication details and financial
entries, while the application returns dashboards, transaction views,
budget status, goal progress and reports. The database stores and retrieves
the corresponding records. Figure~\ref{fig:dfd} summarises this boundary.

\begin{figure}[htbp]
\centering
\begin{center}
\fbox{\begin{minipage}{0.92\textwidth}
\centering
\textbf{CONTEXT-LEVEL DATA-FLOW DIAGRAM}

\vspace{0.25cm}
\begin{tabular}{|c|c|c|}
\hline
\textbf{External entity} & \textbf{Application process} & \textbf{Data store}\\
\hline
\textbf{User} & \textbf{Personal Finance Manager} & \textbf{MySQL}\\
\hline
Login details, income, & $\longrightarrow$ & Read / write\\
expenses, budgets, goals & & queries\\
\hline
Dashboard, reports, & $\longleftarrow$ & Stored records\\
validation messages & &\\
\hline
\end{tabular}
\end{minipage}}
\end{center}
\caption{Context-level data-flow diagram}
\label{fig:dfd}
\end{figure}

\subsection{Database Design Notes}

The physical database, \texttt{personal\_finance\_manager}, contains ten
tables: \texttt{users}, \texttt{categories}, \texttt{income},
\texttt{expenses}, \texttt{transactions}, \texttt{budgets},
\texttt{savings\_goals}, \texttt{goal\_contributions},
\texttt{recurring\_expenses} and \texttt{reports}. The ER model and
schema summary in this chapter were prepared from the supplied MySQL
structure-only definitions, including primary keys, indexes, unique
constraints, data types and declared foreign keys.

The schema stores income and expenses in separate tables, and also
contains a \texttt{transactions} table with a transaction type field.
The \texttt{budgets} table stores a username, category and monthly
limit. Savings targets are stored in \texttt{savings\_goals}, while
individual deposits are represented by \texttt{goal\_contributions}.
The \texttt{recurring\_expenses} table stores recurrence frequency,
next due date and active status. The \texttt{reports} table contains
monthly report metadata and summary totals.

Not every logical association is enforced by a foreign key. The
referential-integrity observations and ER diagram explicitly distinguish
declared foreign keys from application-level associations so that the
report reflects the supplied SQL schema accurately.


\section{Application Interface and Screenshots}
\label{sec:application-screenshots}

This section presents screenshots of the implemented Personal Finance
Manager desktop application. The screenshots document the principal
user interfaces and provide visual evidence of the features discussed
in the implementation and design sections. The image files are expected
to be available in the same directory as \texttt{main.tex}, or inside
the \texttt{figures/} directory.

\subsection{Login Interface}

The login screen is the entry point to the application. It accepts the
registered username and password and provides access to the finance
management modules after successful authentication.

\begin{figure}[htbp]
  \centering
  \includegraphics[width=0.78\textwidth]{login.png}
  \caption{Login interface of the Personal Finance Manager}
  \label{fig:login}
\end{figure}

\subsection{Dashboard}

The dashboard provides an overview of the user's financial position.
It brings together income, expenses, net savings, spending trends,
recent transactions, expense distribution and savings-goal progress.

\begin{figure}[htbp]
  \centering
  \includegraphics[width=\textwidth,height=0.62\textheight,keepaspectratio]{dashboard.png}
  \caption{Dashboard showing financial summaries, charts and recent activity}
  \label{fig:dashboard}
\end{figure}

\subsection{Transaction History}

The transactions interface presents income and expense entries in a
combined view. It helps the user review recorded financial activity
and distinguish transaction types.

\begin{figure}[htbp]
  \centering
  \includegraphics[width=\textwidth,height=0.58\textheight,keepaspectratio]{transactions.png}
  \caption{Transaction history interface}
  \label{fig:transactions}
\end{figure}

\subsection{Income Management}

The income interface allows the user to record income details such as
amount, category, description, date and source. Summary information
helps the user review income over time.

\begin{figure}[htbp]
  \centering
  \includegraphics[width=\textwidth,height=0.58\textheight,keepaspectratio]{income.png}
  \caption{Income management interface}
  \label{fig:income}
\end{figure}

\subsection{Expense Management}

The expense interface supports recording and reviewing expenditure.
The form captures the expense name, amount, category and date, while
the page displays previously recorded expenses.

\begin{figure}[htbp]
  \centering
  \includegraphics[width=\textwidth,height=0.58\textheight,keepaspectratio]{expense.png}
  \caption{Expense management interface}
  \label{fig:expenses}
\end{figure}

\subsection{Budget Management}

The budget interface allows the user to define monthly spending limits
for categories and monitor spending against those limits. It presents
budget usage and highlights cases where spending exceeds a defined limit.

\begin{figure}[htbp]
  \centering
  \includegraphics[width=\textwidth,height=0.58\textheight,keepaspectratio]{budget.png}
  \caption{Budget management interface}
  \label{fig:budget}
\end{figure}

\subsection{Savings Goals}

The savings-goals interface allows the user to define financial targets,
record contributions and monitor progress towards each goal.

\begin{figure}[htbp]
  \centering
  \includegraphics[width=\textwidth,height=0.58\textheight,keepaspectratio]{savings.png}
  \caption{Savings goals and contribution tracking interface}
  \label{fig:savings}
\end{figure}

\subsection{Financial Reports}

The reports interface presents a selected month's income, expenses,
net savings and category-wise expenditure summary. The report can also
be exported to CSV for further use.

\begin{figure}[htbp]
  \centering
  \includegraphics[width=\textwidth,height=0.58\textheight,keepaspectratio]{reports.png}
  \caption{Monthly financial reports interface}
  \label{fig:reports}
\end{figure}

\subsection{Settings}

The settings interface displays account-related information and
provides the available account-management option.

\begin{figure}[htbp]
  \centering
  \includegraphics[width=0.82\textwidth]{settings.png}
  \caption{Application settings interface}
  \label{fig:settings}
\end{figure}

% ===== CHAPTER 3: RESULTS & EVALUATION =====

\chapter{Results and Evaluation}

\section{Testing Methodology}

\begin{itemize}
  \item \textbf{Functional testing:} each page was exercised
    manually using the sample dataset visible in the screenshots
    (income and expense records for October and November 2026,
    three budgets and three savings goals).
  \item \textbf{Input validation testing:} empty fields, non-numeric
    or non-positive amounts and invalid dates are rejected with a
    warning on the forms.
  \item \textbf{Consistency checks:} figures that appear on more
    than one page (totals, percentages, goal progress) were
    compared across pages.
  \item \textbf{Planned:} unit tests for the calculation functions,
    performance measurement of dashboard refresh latency, and a
    pilot with 3--5 peer testers, as defined in the
    proposal~\cite{pfmproposal}.
\end{itemize}

Table~\ref{tab:testcases} lists the manual test cases and their
outcomes.

% Test outcomes below reflect the evaluation recorded during development.
% Re-run the cases on the final build before submission and update any changed result.
\begin{table}[htbp]
\centering
\small
\begin{tabularx}{\textwidth}{|c|l|X|X|c|}
  \hline
  \textbf{ID} & \textbf{Area} & \textbf{Input / action} &
  \textbf{Observed behaviour} & \textbf{Result} \\
  \hline
  T1 & Login & Empty username or password &
    Inline error message shown & Pass \\
  \hline
  T2 & Income form & Amount empty, non-numeric, zero or negative &
    Warning shown; nothing saved & Pass \\
  \hline
  T3 & Income form & Date not in \texttt{dd-mm-yyyy} format &
    Warning shown; nothing saved & Pass \\
  \hline
  T4 & Budget usage & Compare Budget page with sidebar card &
    146\% (Rs.\,19,000) vs 150\% (Rs.\,19,500) & Fail \\
  \hline
  T5 & Goal progress & Compare dashboard with Savings Goals page &
    16\% vs 17\% for the same goal & Fail \\
  \hline
\end{tabularx}
\caption{Manual test cases and outcomes}
\label{tab:testcases}
\end{table}

\section{Implementation Status}

Table~\ref{tab:status} compares the proposal's deliverables with
the current build.

\begin{table}[htbp]
\centering
\begin{tabular}{|l|l|}
  \hline
  \textbf{Deliverable (from proposal)} & \textbf{Status} \\
  \hline
  Login with MySQL authentication & Implemented (plain-text password) \\
  \hline
  Add / view income and expense records & Implemented \\
  \hline
  Dashboard with totals and charts & Implemented \\
  \hline
  Budget engine with per-category limits & Implemented \\
  \hline
  Savings goals with progress & Implemented \\
  \hline
  Reports and CSV export & Implemented \\
  \hline
  User registration & Not yet implemented \\
  \hline
  Recurring expenses & Not yet implemented \\
  \hline
  Password hashing & Not yet implemented \\
  \hline
  Unit tests and latency measurement & Planned \\
  \hline
  Matplotlib charts & Replaced by Tkinter canvas drawing \\
  \hline
\end{tabular}
\caption{Implementation status against the project proposal}
\label{tab:status}
\end{table}

\section{Performance Metrics}

The proposal sets a Dashboard Refresh Latency (DRL) target of
at most one second for up to 5,000 transactions. This has not yet
been measured on the current build; Table~\ref{tab:performance}
will be filled in once timing is added.

% Do not report estimated performance values. Populate this table only
% after running the corresponding measurements on the final build.
\begin{table}[htbp]
\centering
\begin{tabular}{|l|r|r|}
  \hline
  \textbf{Metric} & \textbf{Target} & \textbf{Achieved} \\
  \hline
  Dashboard refresh latency (5,000 records) & $\leq 1$ s & To be measured \\
  \hline
  Calculation accuracy on fixed test set & $100\%$ & To be measured \\
  \hline
\end{tabular}
\caption{Evaluation metrics defined in the proposal}
\label{tab:performance}
\end{table}

\section{Analysis}
\label{sec:analysis}

\begin{itemize}
  \item \textbf{Objectives met:} login (without password hashing),
    transaction entry, dashboard, budgets, savings goals and
    reports are working end to end against MySQL.
  \item \textbf{Challenges:} early versions rebuilt every page on
    each navigation and showed a blank content area when frames
    were embedded incorrectly; moving to a single controller with
    stacked frames and \texttt{on\_show()} refresh hooks fixed this.
  \item \textbf{Inconsistencies found during testing:}
    \begin{itemize}
      \item The Budget page reports 146\% usage (Rs.\,19,000 of
        Rs.\,13,000) while the sidebar spending-limit card reports
        150\% (Rs.\,19,500 of Rs.\,13,000). The Budget page counts
        only categories that have a budget, whereas the sidebar
        counts all expenses, so the two figures differ by the
        Rs.\,500 spent on an unbudgeted category.
      \item Goal percentages differ between pages: the same goal
        (Rs.\,10,000 saved of Rs.\,60,000, i.e.\ 16.67\%) appears as
        16\% on the dashboard and 17\% on the Savings Goals page.
        This is consistent with one view truncating and the other
        rounding to the nearest integer.
      \item Date formats differ between forms: income and expense
        dates use \texttt{dd-mm-yyyy}, while the goal deadline uses
        \texttt{YYYY-MM-DD}.
    \end{itemize}
    The first two should be unified in a single shared calculation
    function and the third in a single date-format setting.
\end{itemize}

\section{Limitations}
\label{sec:limitations}

\begin{itemize}
  \item Passwords are currently compared as plain text and the
    database credentials are set inside the source code. This was
    accepted for the first phase and must be replaced with hashed
    passwords and external configuration before real use.
  \item There is no registration screen and no recurring-expense
    support yet.
  \item All records are entered manually; there is no import from
    bank statements or UPI apps.
  \item Several queries and calculations still live inside the GUI
    classes, so the business-logic tier is not yet fully separate.
  \item Some user/category associations in the schema are not declared as foreign keys (for example, income, expenses and budgets use username fields; reports.user\_id and transactions.category\_id also lack foreign-key constraints).
  \item Dashboard latency and calculation accuracy have not yet been
    measured (Table~\ref{tab:performance}).
\end{itemize}

% ===== CHAPTER 4: CONCLUSION =====

\chapter{Conclusion}

\section{Summary of Work}

A working desktop Personal Finance Manager has been built with
Python, Tkinter and MySQL. It supports login, income and expense
entry, a dashboard, per-category budgets, savings goals, monthly
reports with CSV export, and a settings page. Objectives~1 to~5
are met, although the login does not yet hash passwords;
Objective~6 is partly met.

\section{Key Findings}

\begin{itemize}
  \item \textbf{Finding 1:} building the transaction-and-database
    core first, as planned in the proposal, made it straightforward
    to add budgets, goals and reports on top of the same data.
  \item \textbf{Finding 2:} a single controller with stacked frames
    and refresh hooks is a more reliable navigation pattern in
    Tkinter than recreating widgets on each page change.
  \item \textbf{Finding 3:} the same figure computed in more than
    one place drifts apart (146\% versus 150\% usage), so shared
    calculation functions are needed.
\end{itemize}

\section{Future Work and Recommendations}

\begin{enumerate}
  \item Hash passwords, move credentials to a configuration file,
    and add user registration.
  \item Add recurring expenses with due-date reminders.
  \item Move queries and calculations into a separate data-access
    and logic layer, and add unit tests for the formulas.
  \item Measure dashboard refresh latency on 5,000 records and add
    database indexes on user, date and category columns.
  \item Expose the logic layer through a Flask REST API, following
    the proposal's path from desktop MVP to web, multi-user cloud
    and mobile versions~\cite{pfmproposal}.
\end{enumerate}

\section{Final Remarks}

The project shows that a simple, auditable finance tool can be
built with standard tools and extended in stages. The next phase
is to harden security, remove duplicated calculations and
measure performance against the targets set in the proposal.

% ===== BIBLIOGRAPHY =====

\printbibliography[heading=bibintoc]

\end{document}


